% Tabla de requerimientos funcionales
\begingroup
\renewcommand{\arraystretch}{1.3}
\setlength{\tabcolsep}{6pt}
\centering
\begin{longtable}{|p{0.25\linewidth}|p{0.15\linewidth}|p{0.15\linewidth}|p{0.35\linewidth}|}
\caption{RF01 - AUTENTICACIÓN Y CONTROL DE SESIÓN}\label{tab:rf01}\\
\hline
		\rowcolor{tableheaderblue}
		\multicolumn{4}{|l|}{\textbf{Requerimientos funcionales}}                                                                                                                                           \\ \hline
		\multicolumn{4}{|l|}{\textbf{Prioridad: Alta}}                                                                                                                                                      \\ \hline
		\textbf{CÓDIGO DEL REQUERIMIENTO:}                          & RF01                                                            & \cellcolor{tableheaderblue}\textbf{Actor} & Administrador / Usuario \\ \hline
		\textbf{NOMBRE DEL REQUERIMIENTO}                           & \multicolumn{3}{l|}{\textbf{Autenticación y control de sesión}}                                                                       \\ \hline
		\rowcolor{tableheaderblue}
		\multicolumn{4}{|l|}{\textbf{Descripción}}                                                                                                                                                          \\ \hline
		\multicolumn{4}{|p{0.96\linewidth}|}{%
		El sistema debe gestionar el acceso seguro de los usuarios mediante credenciales y mecanismos de control de sesión basados en tokens.}                                                              \\ \hline
		\rowcolor{tableheaderblue}
		\multicolumn{4}{|l|}{\textbf{Funcionalidad}}                                                                                                                                                        \\ \hline
		\multicolumn{4}{|p{0.96\linewidth}|}{%
		El sistema debe validar las credenciales de acceso (correo y contraseña) contra la base de datos antes de permitir el ingreso del usuario. \newline
		El sistema debe generar un token de sesión (por ejemplo, JWT) para mantener la autenticación del usuario durante su navegación por el sistema. \newline
		El sistema debe restringir el acceso a las rutas internas únicamente a usuarios que posean un token de sesión válido y no expirado.}                                                                \\ \hline
		\cellcolor{tableheaderblue}\textbf{Criterios de aceptación} & \multicolumn{3}{p{0.70\linewidth}|}{                                                                                                  %
			1. El sistema debe permitir que el usuario ingrese su correo y contraseña en un formulario seguro (canal cifrado). \newline
			2. El sistema debe redirigir al dashboard principal cuando las credenciales sean válidas y el token se genere correctamente. \newline
		3. El sistema debe bloquear y redirigir a la pantalla de inicio de sesión cuando se intente acceder a recursos internos sin un token válido.}                                                       \\ \hline
		\rowcolor{tableheaderblue}
		\multicolumn{4}{|l|}{\textbf{Restricciones}}                                                                                                                                                        \\ \hline
		\multicolumn{4}{|p{0.96\linewidth}|}{%
		El sistema debe almacenar las contraseñas encriptadas en la base de datos utilizando algoritmos de hashing seguros.}                                                                                \\ \hline
	
\end{longtable}
\endgroup


\begingroup
\renewcommand{\arraystretch}{1.3}
\setlength{\tabcolsep}{6pt}
\centering
\begin{longtable}{|p{0.25\linewidth}|p{0.15\linewidth}|p{0.15\linewidth}|p{0.35\linewidth}|}
\caption{RF02 - GESTIÓN DEL CATÁLOGO DE PRODUCTOS}\label{tab:rf02}\\
\hline
		\rowcolor{tableheaderblue}
		\multicolumn{4}{|l|}{\textbf{Requerimientos funcionales}}                                                                                                                                    \\ \hline
		\multicolumn{4}{|l|}{\textbf{Prioridad: Alta}}                                                                                                                                               \\ \hline
		\textbf{CÓDIGO DEL REQUERIMIENTO:}                          & RF02                                                            & \cellcolor{tableheaderblue}\textbf{Actor} & Administrador    \\ \hline
		\textbf{NOMBRE DEL REQUERIMIENTO}                           & \multicolumn{3}{l|}{\textbf{Gestión del catálogo de productos}}                                                                \\ \hline
		\rowcolor{tableheaderblue}
		\multicolumn{4}{|l|}{\textbf{Descripción}}                                                                                                                                                   \\ \hline
		\multicolumn{4}{|p{0.96\linewidth}|}{%
		El sistema debe permitir la administración completa del catálogo de productos comercializados, garantizando la unicidad de identificadores y la disponibilidad inmediata de la información.} \\ \hline
		\rowcolor{tableheaderblue}
		\multicolumn{4}{|l|}{\textbf{Funcionalidad}}                                                                                                                                                 \\ \hline
		\multicolumn{4}{|p{0.96\linewidth}|}{%
		El sistema debe permitir crear, leer, actualizar y desactivar productos, capturando como mínimo: nombre, SKU, código de barras, categoría, precio de costo, precio de venta, impuesto aplicable y stock actual. \newline
		El sistema debe validar que el SKU y el código de barras de un producto sean únicos dentro del catálogo. \newline
		El sistema debe aplicar eliminación lógica (borrado lógico) a los productos con ventas históricas, evitando la pérdida de información asociada.}                                             \\ \hline
		\cellcolor{tableheaderblue}\textbf{Criterios de aceptación} & \multicolumn{3}{p{0.70\linewidth}|}{                                                                                           %
			1. El sistema debe mostrar un mensaje de error cuando se intente registrar un producto con un SKU o código de barras ya existente. \newline
			2. El sistema debe mostrar el producto recién creado o actualizado de forma inmediata en el listado general de inventario. \newline
		3. El sistema debe impedir la eliminación física de productos que tengan ventas asociadas, permitiendo únicamente su desactivación.}                                                         \\ \hline
		\rowcolor{tableheaderblue}
		\multicolumn{4}{|l|}{\textbf{Restricciones}}                                                                                                                                                 \\ \hline
		\multicolumn{4}{|p{0.96\linewidth}|}{%
		El sistema debe manejar paginación en el listado de productos cuando el número de registros supere el límite establecido.}                                                                   \\ \hline
	
\end{longtable}
\endgroup


\begingroup
\renewcommand{\arraystretch}{1.3}
\setlength{\tabcolsep}{6pt}
\centering
\begin{longtable}{|p{0.25\linewidth}|p{0.15\linewidth}|p{0.15\linewidth}|p{0.35\linewidth}|}
\caption{RF03 - REGISTRO DE MOVIMIENTOS DE INVENTARIO}\label{tab:rf03}\\
\hline
		\rowcolor{tableheaderblue}
		\multicolumn{4}{|l|}{\textbf{Requerimientos funcionales}}                                                                                                                                                             \\ \hline
		\multicolumn{4}{|l|}{\textbf{Prioridad: Alta}}                                                                                                                                                                        \\ \hline
		\textbf{CÓDIGO DEL REQUERIMIENTO:}                          & RF03                                                                & \cellcolor{tableheaderblue}\textbf{Actor} & Administrador / Usuario de inventario \\ \hline
		\textbf{NOMBRE DEL REQUERIMIENTO}                           & \multicolumn{3}{l|}{\textbf{Registro de movimientos de inventario}}                                                                                     \\ \hline
		\rowcolor{tableheaderblue}
		\multicolumn{4}{|l|}{\textbf{Descripción}}                                                                                                                                                                            \\ \hline
		\multicolumn{4}{|p{0.96\linewidth}|}{%
		El sistema debe mantener la trazabilidad de las variaciones de inventario distintas a las ventas, mediante el registro estructurado de movimientos de entrada y salida.}                                              \\ \hline
		\rowcolor{tableheaderblue}
		\multicolumn{4}{|l|}{\textbf{Funcionalidad}}                                                                                                                                                                          \\ \hline
		\multicolumn{4}{|p{0.96\linewidth}|}{%
		El sistema debe registrar entradas y salidas de inventario clasificándolas por tipo (compra, ajuste, daño, devolución u otros definidos). \newline
		El sistema debe actualizar en tiempo cercano al real el stock disponible del producto según la cantidad registrada en cada movimiento. \newline
		El sistema debe almacenar para cada movimiento la fecha, la hora, el tipo de movimiento, el usuario responsable y una justificación.}                                                                                 \\ \hline
		\cellcolor{tableheaderblue}\textbf{Criterios de aceptación} & \multicolumn{3}{p{0.70\linewidth}|}{                                                                                                                    %
			1. El sistema debe impedir registrar movimientos de salida que dejen el stock resultante en un valor negativo. \newline
			2. El sistema debe reflejar inmediatamente en la ficha del producto el nuevo stock disponible tras registrar un movimiento. \newline
		3. El sistema debe permitir consultar el historial de movimientos por producto y por rango de fechas.}                                                                                                                \\ \hline
		\rowcolor{tableheaderblue}
		\multicolumn{4}{|l|}{\textbf{Restricciones}}                                                                                                                                                                          \\ \hline
		\multicolumn{4}{|p{0.96\linewidth}|}{%
		El sistema debe tratar cada movimiento como irreversible una vez confirmado; para corregir errores se debe generar un contra-movimiento con signo contrario.}                                                         \\ \hline
	
\end{longtable}
\endgroup


\begingroup
\renewcommand{\arraystretch}{1.3}
\setlength{\tabcolsep}{6pt}
\centering
\begin{longtable}{|p{0.25\linewidth}|p{0.15\linewidth}|p{0.15\linewidth}|p{0.35\linewidth}|}
\caption{RF04 - ALERTAS DE STOCK BAJO Y UMBRAL MÍNIMO}\label{tab:rf04}\\
\hline
		\rowcolor{tableheaderblue}
		\multicolumn{4}{|l|}{\textbf{Requerimientos funcionales}}                                                                                                                                                                       \\ \hline
		\multicolumn{4}{|l|}{\textbf{Prioridad: Media}}                                                                                                                                                                                 \\ \hline
		\textbf{CÓDIGO DEL REQUERIMIENTO:}                          & RF04                                                                & \cellcolor{tableheaderblue}\textbf{Actor} & Sistema / Administrador / Usuario de inventario \\ \hline
		\textbf{NOMBRE DEL REQUERIMIENTO}                           & \multicolumn{3}{l|}{\textbf{Alertas de stock bajo y umbral mínimo}}                                                                                               \\ \hline
		\rowcolor{tableheaderblue}
		\multicolumn{4}{|l|}{\textbf{Descripción}}                                                                                                                                                                                      \\ \hline
		\multicolumn{4}{|p{0.96\linewidth}|}{%
		El sistema debe notificar de forma proactiva sobre productos cuyo nivel de existencias se encuentre por debajo del umbral mínimo configurado, facilitando la toma de decisiones de reabastecimiento.}                           \\ \hline
		\rowcolor{tableheaderblue}
		\multicolumn{4}{|l|}{\textbf{Funcionalidad}}                                                                                                                                                                                    \\ \hline
		\multicolumn{4}{|p{0.96\linewidth}|}{%
		El sistema debe permitir configurar un umbral mínimo de stock por producto. \newline
		El sistema debe identificar y marcar como “stock bajo” los productos cuyo stock actual sea inferior o igual al umbral configurado. \newline
		El sistema debe resaltar visualmente los productos con stock bajo en el listado de inventario y/o en el dashboard principal.}                                                                                                   \\ \hline
		\cellcolor{tableheaderblue}\textbf{Criterios de aceptación} & \multicolumn{3}{p{0.70\linewidth}|}{                                                                                                                              %
			1. El sistema debe recalcular las alertas de stock automáticamente cada vez que se registre una venta o un movimiento de inventario. \newline
			2. El sistema debe resaltar claramente los productos en condición de stock bajo mediante iconos, colores o etiquetas diferenciadas. \newline
		3. El sistema debe permitir que el usuario modifique el umbral mínimo de un producto y que la alerta se actualice en consecuencia.}                                                                                             \\ \hline
		\rowcolor{tableheaderblue}
		\multicolumn{4}{|l|}{\textbf{Restricciones}}                                                                                                                                                                                    \\ \hline
		\multicolumn{4}{|p{0.96\linewidth}|}{%
		El sistema no debe generar alertas para productos marcados como inactivos.}                                                                                                                                                     \\ \hline
	
\end{longtable}
\endgroup


\begingroup
\renewcommand{\arraystretch}{1.3}
\setlength{\tabcolsep}{6pt}
\centering
\begin{longtable}{|p{0.25\linewidth}|p{0.15\linewidth}|p{0.15\linewidth}|p{0.35\linewidth}|}
\caption{RF05 - PROCESAMIENTO DE VENTA EN TIEMPO REAL}\label{tab:rf05}\\
\hline
		\rowcolor{tableheaderblue}
		\multicolumn{4}{|l|}{\textbf{Requerimientos funcionales}}                                                                                                                                                                    \\ \hline
		\multicolumn{4}{|l|}{\textbf{Prioridad: Alta}}                                                                                                                                                                               \\ \hline
		\textbf{CÓDIGO DEL REQUERIMIENTO:}                          & RF05                                                                & \cellcolor{tableheaderblue}\textbf{Actor} & Cajero / Administrador                       \\ \hline
		\textbf{NOMBRE DEL REQUERIMIENTO}                           & \multicolumn{3}{l|}{\textbf{Procesamiento de venta en tiempo real}}                                                                                            \\ \hline
		\rowcolor{tableheaderblue}
		\multicolumn{4}{|l|}{\textbf{Descripción}}                                                                                                                                                                                   \\ \hline
		\multicolumn{4}{|p{0.96\linewidth}|}{%
		El sistema debe proporcionar una interfaz ágil para la ejecución de ventas, permitiendo la selección de productos, el cálculo automático de totales y el registro de la transacción con impacto inmediato en el inventario.} \\ \hline
		\rowcolor{tableheaderblue}
		\multicolumn{4}{|l|}{\textbf{Funcionalidad}}                                                                                                                                                                                 \\ \hline
		\multicolumn{4}{|p{0.96\linewidth}|}{%
		El sistema debe permitir buscar productos por nombre, SKU o código de barras y agregarlos a un carrito de venta con la cantidad correspondiente. \newline
		El sistema debe calcular automáticamente el subtotal, el valor del impuesto y el total final de la venta según los datos del catálogo. \newline
		El sistema debe registrar la venta con un identificador único, fecha, hora, usuario, productos vendidos, cantidades, totales y cliente asociado (cuando aplique).}                                                           \\ \hline
		\cellcolor{tableheaderblue}\textbf{Criterios de aceptación} & \multicolumn{3}{p{0.70\linewidth}|}{                                                                                                                           %
			1. El sistema debe impedir confirmar una venta cuando al menos uno de los productos no disponga de stock suficiente. \newline
			2. El sistema debe mostrar al usuario un resumen detallado de la venta antes de su confirmación. \newline
		3. El sistema debe completar el registro de la venta y dejarla disponible para consulta en el historial de ventas inmediatamente después de su confirmación.}                                                                \\ \hline
		\rowcolor{tableheaderblue}
		\multicolumn{4}{|l|}{\textbf{Restricciones}}                                                                                                                                                                                 \\ \hline
		\multicolumn{4}{|p{0.96\linewidth}|}{%
		La interfaz de venta debe minimizar el número de clics necesarios para completar una transacción típica en el punto de venta.}                                                                                               \\ \hline
	
\end{longtable}
\endgroup


\begingroup
\renewcommand{\arraystretch}{1.3}
\setlength{\tabcolsep}{6pt}
\centering
\begin{longtable}{|p{0.25\linewidth}|p{0.15\linewidth}|p{0.15\linewidth}|p{0.35\linewidth}|}
\caption{RF06 - GESTIÓN DE MÉTODOS DE PAGO Y CÁLCULO DE CAMBIO}\label{tab:rf06}\\
\hline
		\rowcolor{tableheaderblue}
		\multicolumn{4}{|l|}{\textbf{Requerimientos funcionales}}                                                                                                                                                       \\ \hline
		\multicolumn{4}{|l|}{\textbf{Prioridad: Alta}}                                                                                                                                                                  \\ \hline
		\textbf{CÓDIGO DEL REQUERIMIENTO:}                          & RF06                                                                         & \cellcolor{tableheaderblue}\textbf{Actor} & Cajero / Administrador \\ \hline
		\textbf{NOMBRE DEL REQUERIMIENTO}                           & \multicolumn{3}{l|}{\textbf{Gestión de métodos de pago y cálculo de cambio}}                                                                      \\ \hline
		\rowcolor{tableheaderblue}
		\multicolumn{4}{|l|}{\textbf{Descripción}}                                                                                                                                                                      \\ \hline
		\multicolumn{4}{|p{0.96\linewidth}|}{%
		El sistema debe manejar diferentes métodos de pago y calcular el cambio a devolver al cliente cuando el pago se realice en efectivo.}                                                                           \\ \hline
		\rowcolor{tableheaderblue}
		\multicolumn{4}{|l|}{\textbf{Funcionalidad}}                                                                                                                                                                    \\ \hline
		\multicolumn{4}{|p{0.96\linewidth}|}{%
		El sistema debe permitir seleccionar el método de pago para cada venta (efectivo, tarjeta, transferencia u otros definidos). \newline
		El sistema debe permitir ingresar el monto entregado por el cliente en caso de pago en efectivo y calcular automáticamente el cambio a devolver. \newline
		El sistema debe registrar el método de pago y los valores asociados en el detalle de la venta para su posterior análisis.}                                                                                      \\ \hline
		\cellcolor{tableheaderblue}\textbf{Criterios de aceptación} & \multicolumn{3}{p{0.70\linewidth}|}{                                                                                                              %
			1. El sistema debe validar que el monto entregado por el cliente en efectivo sea igual o superior al total de la venta. \newline
			2. El sistema debe mostrar de forma clara el valor del cambio a devolver antes de finalizar la transacción. \newline
		3. El sistema debe permitir filtrar ventas por método de pago en los reportes y consultas.}                                                                                                                     \\ \hline
		\rowcolor{tableheaderblue}
		\multicolumn{4}{|l|}{\textbf{Restricciones}}                                                                                                                                                                    \\ \hline
		\multicolumn{4}{|p{0.96\linewidth}|}{%
		Los métodos de pago disponibles deben ser configurables únicamente por usuarios con permisos de administración.}                                                                                                \\ \hline
	
\end{longtable}
\endgroup


\begingroup
\renewcommand{\arraystretch}{1.3}
\setlength{\tabcolsep}{6pt}
\centering
\begin{longtable}{|p{0.25\linewidth}|p{0.15\linewidth}|p{0.15\linewidth}|p{0.35\linewidth}|}
\caption{RF07 - ACTUALIZACIÓN AUTOMÁTICA DE STOCK POR VENTA}\label{tab:rf07}\\
\hline
		\rowcolor{tableheaderblue}
		\multicolumn{4}{|l|}{\textbf{Requerimientos funcionales}}                                                                                                                                                              \\ \hline
		\multicolumn{4}{|l|}{\textbf{Prioridad: Alta}}                                                                                                                                                                         \\ \hline
		\textbf{CÓDIGO DEL REQUERIMIENTO:}                          & RF07                                                                      & \cellcolor{tableheaderblue}\textbf{Actor} & Sistema / Cajero / Administrador \\ \hline
		\textbf{NOMBRE DEL REQUERIMIENTO}                           & \multicolumn{3}{l|}{\textbf{Actualización automática de stock por venta}}                                                                                \\ \hline
		\rowcolor{tableheaderblue}
		\multicolumn{4}{|l|}{\textbf{Descripción}}                                                                                                                                                                             \\ \hline
		\multicolumn{4}{|p{0.96\linewidth}|}{%
		El sistema debe actualizar automáticamente las existencias de los productos involucrados en una venta al momento de confirmar la transacción, manteniendo la consistencia con el módulo de inventario.}                \\ \hline
		\rowcolor{tableheaderblue}
		\multicolumn{4}{|l|}{\textbf{Funcionalidad}}                                                                                                                                                                           \\ \hline
		\multicolumn{4}{|p{0.96\linewidth}|}{%
		El sistema debe descontar del stock de cada producto la cantidad vendida cuando una venta se confirme exitosamente. \newline
		El sistema debe registrar un movimiento de inventario de tipo “salida por venta” asociado al identificador de la venta. \newline
		El sistema debe impedir la confirmación de la venta si, por concurrencia u otros factores, el stock disponible es insuficiente en el momento del cierre.}                                                              \\ \hline
		\cellcolor{tableheaderblue}\textbf{Criterios de aceptación} & \multicolumn{3}{p{0.70\linewidth}|}{                                                                                                                     %
			1. El sistema debe mostrar el nuevo stock disponible del producto inmediatamente después de confirmar la venta. \newline
			2. El sistema debe impedir la confirmación de la venta si no se logra aplicar correctamente el descuento de inventario. \newline
		3. El sistema debe permitir consultar, desde la ficha de un producto, los movimientos de salida generados por ventas.}                                                                                                 \\ \hline
		\rowcolor{tableheaderblue}
		\multicolumn{4}{|l|}{\textbf{Restricciones}}                                                                                                                                                                           \\ \hline
		\multicolumn{4}{|p{0.96\linewidth}|}{%
		La operación de registro de venta y descuento de inventario debe ejecutarse de forma atómica para evitar inconsistencias.}                                                                                             \\ \hline
	
\end{longtable}
\endgroup


\begingroup
\renewcommand{\arraystretch}{1.3}
\setlength{\tabcolsep}{6pt}
\centering
\begin{longtable}{|p{0.25\linewidth}|p{0.15\linewidth}|p{0.15\linewidth}|p{0.35\linewidth}|}
\caption{RF08 - EMISIÓN DE COMPROBANTES DE VENTA}\label{tab:rf08}\\
\hline
		\rowcolor{tableheaderblue}
		\multicolumn{4}{|l|}{\textbf{Requerimientos funcionales}}                                                                                                                                         \\ \hline
		\multicolumn{4}{|l|}{\textbf{Prioridad: Media}}                                                                                                                                                   \\ \hline
		\textbf{CÓDIGO DEL REQUERIMIENTO:}                          & RF08                                                           & \cellcolor{tableheaderblue}\textbf{Actor} & Cajero / Administrador \\ \hline
		\textbf{NOMBRE DEL REQUERIMIENTO}                           & \multicolumn{3}{l|}{\textbf{Emisión de comprobantes de venta}}                                                                      \\ \hline
		\rowcolor{tableheaderblue}
		\multicolumn{4}{|l|}{\textbf{Descripción}}                                                                                                                                                        \\ \hline
		\multicolumn{4}{|p{0.96\linewidth}|}{%
		El sistema debe generar un comprobante documental de cada venta realizada, disponible en formato imprimible o digital, para soporte del cliente y del negocio.}                                   \\ \hline
		\rowcolor{tableheaderblue}
		\multicolumn{4}{|l|}{\textbf{Funcionalidad}}                                                                                                                                                      \\ \hline
		\multicolumn{4}{|p{0.96\linewidth}|}{%
		El sistema debe generar automáticamente un comprobante al finalizar exitosamente una venta, incluyendo: logo del negocio, fecha, número de venta, datos del cliente (si aplica), detalle de ítems, impuestos, total y método de pago. \newline
		El sistema debe permitir la visualización del comprobante desde el navegador y su impresión directa o descarga en formato digital (por ejemplo, PDF). \newline
		El sistema debe manejar un consecutivo incremental y único para la numeración de las ventas.}                                                                                                     \\ \hline
		\cellcolor{tableheaderblue}\textbf{Criterios de aceptación} & \multicolumn{3}{p{0.70\linewidth}|}{                                                                                                %
			1. El sistema debe generar el comprobante inmediatamente después de finalizar la venta sin requerir acciones adicionales del usuario. \newline
			2. El sistema debe presentar el comprobante en un formato compatible con impresoras de tirilla y/o tamaño carta, según configuración. \newline
		3. El sistema debe permitir recuperar y volver a imprimir o descargar el comprobante desde el historial de ventas.}                                                                               \\ \hline
		\rowcolor{tableheaderblue}
		\multicolumn{4}{|l|}{\textbf{Restricciones}}                                                                                                                                                      \\ \hline
		\multicolumn{4}{|p{0.96\linewidth}|}{%
		El diseño del comprobante debe ajustarse a los formatos definidos por la configuración del negocio y las normas vigentes aplicables.}                                                             \\ \hline
	
\end{longtable}
\endgroup


\begingroup
\renewcommand{\arraystretch}{1.3}
\setlength{\tabcolsep}{6pt}
\centering
\begin{longtable}{|p{0.25\linewidth}|p{0.15\linewidth}|p{0.15\linewidth}|p{0.35\linewidth}|}
\caption{RF09 - GESTIÓN DE ANULACIONES Y DEVOLUCIONES}\label{tab:rf09}\\
\hline
		\rowcolor{tableheaderblue}
		\multicolumn{4}{|l|}{\textbf{Requerimientos funcionales}}                                                                                                                                     \\ \hline
		\multicolumn{4}{|l|}{\textbf{Prioridad: Media}}                                                                                                                                               \\ \hline
		\textbf{CÓDIGO DEL REQUERIMIENTO:}                          & RF09                                                                & \cellcolor{tableheaderblue}\textbf{Actor} & Administrador \\ \hline
		\textbf{NOMBRE DEL REQUERIMIENTO}                           & \multicolumn{3}{l|}{\textbf{Gestión de anulaciones y devoluciones}}                                                             \\ \hline
		\rowcolor{tableheaderblue}
		\multicolumn{4}{|l|}{\textbf{Descripción}}                                                                                                                                                    \\ \hline
		\multicolumn{4}{|p{0.96\linewidth}|}{%
		El sistema debe permitir gestionar anulaciones totales y devoluciones parciales de ventas, ajustando en consecuencia el inventario y los registros históricos.}                               \\ \hline
		\rowcolor{tableheaderblue}
		\multicolumn{4}{|l|}{\textbf{Funcionalidad}}                                                                                                                                                  \\ \hline
		\multicolumn{4}{|p{0.96\linewidth}|}{%
		El sistema debe permitir seleccionar una venta registrada y marcarla como anulada o registrar la devolución de uno o varios productos. \newline
		El sistema debe generar los movimientos de inventario necesarios para revertir total o parcialmente el descuento de stock asociado a la venta original. \newline
		El sistema debe registrar la fecha, el usuario responsable y el motivo de la anulación o devolución.}                                                                                         \\ \hline
		\cellcolor{tableheaderblue}\textbf{Criterios de aceptación} & \multicolumn{3}{p{0.70\linewidth}|}{                                                                                            %
			1. El sistema debe reflejar el ajuste del inventario derivado de la anulación o devolución en el stock de los productos afectados. \newline
			2. El sistema debe mantener visible en el historial la venta original con un estado que indique “anulada” o “con devolución”. \newline
		3. El sistema debe impedir la anulación de ventas cuando existan restricciones adicionales configuradas (por ejemplo, cierres de caja o cierres contables).}                                  \\ \hline
		\rowcolor{tableheaderblue}
		\multicolumn{4}{|l|}{\textbf{Restricciones}}                                                                                                                                                  \\ \hline
		\multicolumn{4}{|p{0.96\linewidth}|}{%
		Solo usuarios con permisos de administración pueden ejecutar procesos de anulación o devolución.}                                                                                             \\ \hline
	
\end{longtable}
\endgroup


\begingroup
\renewcommand{\arraystretch}{1.3}
\setlength{\tabcolsep}{6pt}
\centering
\begin{longtable}{|p{0.25\linewidth}|p{0.15\linewidth}|p{0.15\linewidth}|p{0.35\linewidth}|}
\caption{RF10 - GESTIÓN DEL DIRECTORIO DE CLIENTES}\label{tab:rf10}\\
\hline
		\rowcolor{tableheaderblue}
		\multicolumn{4}{|l|}{\textbf{Requerimientos funcionales}}                                                                                                                                           \\ \hline
		\multicolumn{4}{|l|}{\textbf{Prioridad: Media}}                                                                                                                                                     \\ \hline
		\textbf{CÓDIGO DEL REQUERIMIENTO:}                          & RF10                                                             & \cellcolor{tableheaderblue}\textbf{Actor} & Administrador / Cajero \\ \hline
		\textbf{NOMBRE DEL REQUERIMIENTO}                           & \multicolumn{3}{l|}{\textbf{Gestión del directorio de clientes}}                                                                      \\ \hline
		\rowcolor{tableheaderblue}
		\multicolumn{4}{|l|}{\textbf{Descripción}}                                                                                                                                                          \\ \hline
		\multicolumn{4}{|p{0.96\linewidth}|}{%
		El sistema debe centralizar la información de los clientes para facilitar su uso en las ventas y en los procesos de seguimiento comercial.}                                                         \\ \hline
		\rowcolor{tableheaderblue}
		\multicolumn{4}{|l|}{\textbf{Funcionalidad}}                                                                                                                                                        \\ \hline
		\multicolumn{4}{|p{0.96\linewidth}|}{%
		El sistema debe permitir registrar, editar y buscar clientes, almacenando como mínimo: nombre o razón social, documento de identidad, teléfono, correo electrónico y dirección. \newline
		El sistema debe validar que no se registren clientes duplicados con el mismo número de documento de identidad. \newline
		El sistema debe permitir asociar un cliente existente a una venta en el módulo POS y aplicar eliminación lógica cuando se desee inactivar un cliente.}                                              \\ \hline
		\cellcolor{tableheaderblue}\textbf{Criterios de aceptación} & \multicolumn{3}{p{0.70\linewidth}|}{                                                                                                  %
			1. El sistema debe mostrar un mensaje de error cuando se intente registrar un cliente con un número de documento ya existente. \newline
			2. El sistema debe permitir filtrar rápidamente clientes por nombre, documento o correo electrónico. \newline
		3. El sistema debe conservar el historial de ventas asociadas al cliente aun cuando este se marque como inactivo.}                                                                                  \\ \hline
		\rowcolor{tableheaderblue}
		\multicolumn{4}{|l|}{\textbf{Restricciones}}                                                                                                                                                        \\ \hline
		\multicolumn{4}{|p{0.96\linewidth}|}{%
		El tratamiento de la información de clientes debe cumplir con las políticas de protección de datos personales aplicables al negocio.}                                                               \\ \hline
	
\end{longtable}
\endgroup


\begingroup
\renewcommand{\arraystretch}{1.3}
\setlength{\tabcolsep}{6pt}
\centering
\begin{longtable}{|p{0.25\linewidth}|p{0.15\linewidth}|p{0.15\linewidth}|p{0.35\linewidth}|}
\caption{RF11 - CONSULTA DE HISTORIAL DE COMPRAS POR CLIENTE}\label{tab:rf11}\\
\hline
		\rowcolor{tableheaderblue}
		\multicolumn{4}{|l|}{\textbf{Requerimientos funcionales}}                                                                                                                                                     \\ \hline
		\multicolumn{4}{|l|}{\textbf{Prioridad: Media}}                                                                                                                                                               \\ \hline
		\textbf{CÓDIGO DEL REQUERIMIENTO:}                          & RF11                                                                       & \cellcolor{tableheaderblue}\textbf{Actor} & Administrador / Cajero \\ \hline
		\textbf{NOMBRE DEL REQUERIMIENTO}                           & \multicolumn{3}{l|}{\textbf{Consulta de historial de compras por cliente}}                                                                      \\ \hline
		\rowcolor{tableheaderblue}
		\multicolumn{4}{|l|}{\textbf{Descripción}}                                                                                                                                                                    \\ \hline
		\multicolumn{4}{|p{0.96\linewidth}|}{%
		El sistema debe permitir consultar el historial de compras de un cliente específico, con el fin de apoyar el análisis de comportamiento y la fidelización.}                                                   \\ \hline
		\rowcolor{tableheaderblue}
		\multicolumn{4}{|l|}{\textbf{Funcionalidad}}                                                                                                                                                                  \\ \hline
		\multicolumn{4}{|p{0.96\linewidth}|}{%
		El sistema debe permitir seleccionar un cliente y listar las ventas asociadas, mostrando fecha, número de venta, monto total y método de pago. \newline
		El sistema debe permitir filtrar el historial de compras por rango de fechas. \newline
		El sistema debe calcular el total acumulado de compras del cliente en el período consultado.}                                                                                                                 \\ \hline
		\cellcolor{tableheaderblue}\textbf{Criterios de aceptación} & \multicolumn{3}{p{0.70\linewidth}|}{                                                                                                            %
			1. El sistema debe mostrar únicamente las ventas correspondientes al cliente seleccionado y al período filtrado. \newline
			2. El sistema debe actualizar el historial de compras en tiempo cercano al real cuando se registre una nueva venta del cliente. \newline
		3. El sistema debe permitir exportar o imprimir el historial de compras si la funcionalidad de exportación está habilitada.}                                                                                  \\ \hline
		\rowcolor{tableheaderblue}
		\multicolumn{4}{|l|}{\textbf{Restricciones}}                                                                                                                                                                  \\ \hline
		\multicolumn{4}{|p{0.96\linewidth}|}{%
		El acceso al historial detallado de compras debe estar restringido a usuarios con permisos autorizados.}                                                                                                      \\ \hline
	
\end{longtable}
\endgroup


\begingroup
\renewcommand{\arraystretch}{1.3}
\setlength{\tabcolsep}{6pt}
\centering
\begin{longtable}{|p{0.25\linewidth}|p{0.15\linewidth}|p{0.15\linewidth}|p{0.35\linewidth}|}
\caption{RF12 - DASHBOARD DE INDICADORES Y REPORTE HISTÓRICO FILTRABLE}\label{tab:rf12}\\
\hline
		\rowcolor{tableheaderblue}
		\multicolumn{4}{|l|}{\textbf{Requerimientos funcionales}}                                                                                                                                                      \\ \hline
		\multicolumn{4}{|l|}{\textbf{Prioridad: Media}}                                                                                                                                                                \\ \hline
		\textbf{CÓDIGO DEL REQUERIMIENTO:}                          & RF12                                                                                 & \cellcolor{tableheaderblue}\textbf{Actor} & Administrador \\ \hline
		\textbf{NOMBRE DEL REQUERIMIENTO}                           & \multicolumn{3}{l|}{\textbf{Dashboard de indicadores y reporte histórico filtrable}}                                                             \\ \hline
		\rowcolor{tableheaderblue}
		\multicolumn{4}{|l|}{\textbf{Descripción}}                                                                                                                                                                     \\ \hline
		\multicolumn{4}{|p{0.96\linewidth}|}{%
		El sistema debe ofrecer una vista ejecutiva del estado del negocio mediante indicadores clave y reportes históricos de ventas con filtros avanzados.}                                                          \\ \hline
		\rowcolor{tableheaderblue}
		\multicolumn{4}{|l|}{\textbf{Funcionalidad}}                                                                                                                                                                   \\ \hline
		\multicolumn{4}{|p{0.96\linewidth}|}{%
		El sistema debe calcular y mostrar en el dashboard métricas clave como ventas del día, ventas del mes, número de transacciones y número de productos con stock bajo. \newline
		El sistema debe permitir generar reportes históricos de ventas filtrables por rango de fechas, cliente, producto y método de pago, presentando el detalle de cada transacción. \newline
		El sistema debe recalcular los totales monetarios del reporte según los filtros aplicados y presentar la información en tablas paginadas y/o gráficas.}                                                        \\ \hline
		\cellcolor{tableheaderblue}\textbf{Criterios de aceptación} & \multicolumn{3}{p{0.70\linewidth}|}{                                                                                                             %
			1. El sistema debe actualizar los indicadores del dashboard en tiempo cercano al real o al recargar la página. \newline
			2. El sistema debe permitir al usuario seleccionar un rango de fechas y aplicar filtros combinados para la generación del reporte histórico de ventas. \newline
		3. El sistema debe garantizar que los totales mostrados en el dashboard y en los reportes coincidan con la sumatoria de las ventas almacenadas en la base de datos bajo los filtros activos.}                  \\ \hline
		\rowcolor{tableheaderblue}
		\multicolumn{4}{|l|}{\textbf{Restricciones}}                                                                                                                                                                   \\ \hline
		\multicolumn{4}{|p{0.96\linewidth}|}{%
		Las consultas del dashboard y de los reportes deben estar optimizadas para no degradar significativamente el tiempo de respuesta del sistema, aun con volúmenes crecientes de datos.}                          \\ \hline
	
\end{longtable}
\endgroup


\begingroup
\renewcommand{\arraystretch}{1.3}
\setlength{\tabcolsep}{6pt}
\centering
\begin{longtable}{|p{0.25\linewidth}|p{0.15\linewidth}|p{0.15\linewidth}|p{0.35\linewidth}|}
\caption{RF13 - GESTIÓN DE CATEGORÍAS DE PRODUCTOS}\label{tab:rf13}\\
\hline
		\rowcolor{tableheaderblue}
		\multicolumn{4}{|l|}{\textbf{Requerimientos funcionales}}                                                                                                                                         \\ \hline
		\multicolumn{4}{|l|}{\textbf{Prioridad: Media}}                                                                                                                                                   \\ \hline
		\textbf{CÓDIGO DEL REQUERIMIENTO:}                          & RF13                                                             & \cellcolor{tableheaderblue}\textbf{Actor} & Administrador        \\ \hline
		\textbf{NOMBRE DEL REQUERIMIENTO}                           & \multicolumn{3}{l|}{\textbf{Gestión de categorías de productos}}                                                                    \\ \hline
		\rowcolor{tableheaderblue}
		\multicolumn{4}{|l|}{\textbf{Descripción}}                                                                                                                                                        \\ \hline
		\multicolumn{4}{|p{0.96\linewidth}|}{%
		El sistema debe permitir la creación, modificación, consulta y eliminación de categorías para organizar el catálogo de productos, facilitando la navegación y búsqueda por agrupaciones lógicas.} \\ \hline
		\rowcolor{tableheaderblue}
		\multicolumn{4}{|l|}{\textbf{Funcionalidad}}                                                                                                                                                      \\ \hline
		\multicolumn{4}{|p{0.96\linewidth}|}{%
		El sistema debe permitir crear categorías con nombre único y descripción opcional. \newline
		El sistema debe permitir editar el nombre y descripción de categorías existentes. \newline
		El sistema debe permitir desactivar categorías que contengan productos asociados mediante eliminación lógica. \newline
		El sistema debe mostrar la cantidad de productos asociados a cada categoría en el listado de categorías.}                                                                                         \\ \hline
		\cellcolor{tableheaderblue}\textbf{Criterios de aceptación} & \multicolumn{3}{p{0.70\linewidth}|}{                                                                                                %
			1. El sistema debe impedir la creación de categorías con nombres duplicados. \newline
			2. El sistema debe reasignar automáticamente los productos a una categoría "Sin categoría" cuando se elimine lógicamente su categoría original. \newline
		3. El sistema debe permitir filtrar el catálogo de productos por categoría seleccionada.}                                                                                                         \\ \hline
		\rowcolor{tableheaderblue}
		\multicolumn{4}{|l|}{\textbf{Restricciones}}                                                                                                                                                      \\ \hline
		\multicolumn{4}{|p{0.96\linewidth}|}{%
		El sistema debe validar que el nombre de la categoría no exceda los 100 caracteres y que no contenga caracteres especiales no permitidos.}                                                        \\ \hline
	
\end{longtable}
\endgroup


\begingroup
\renewcommand{\arraystretch}{1.3}
\setlength{\tabcolsep}{6pt}
\centering
\begin{longtable}{|p{0.25\linewidth}|p{0.15\linewidth}|p{0.15\linewidth}|p{0.35\linewidth}|}
\caption{RF14 - GESTIÓN DE IMÁGENES DE PRODUCTOS}\label{tab:rf14}\\
\hline
		\rowcolor{tableheaderblue}
		\multicolumn{4}{|l|}{\textbf{Requerimientos funcionales}}                                                                                                                                      \\ \hline
		\multicolumn{4}{|l|}{\textbf{Prioridad: Media}}                                                                                                                                                \\ \hline
		\textbf{CÓDIGO DEL REQUERIMIENTO:}                          & RF14                                                           & \cellcolor{tableheaderblue}\textbf{Actor} & Administrador       \\ \hline
		\textbf{NOMBRE DEL REQUERIMIENTO}                           & \multicolumn{3}{l|}{\textbf{Gestión de imágenes de productos}}                                                                   \\ \hline
		\rowcolor{tableheaderblue}
		\multicolumn{4}{|l|}{\textbf{Descripción}}                                                                                                                                                     \\ \hline
		\multicolumn{4}{|p{0.96\linewidth}|}{%
		El sistema debe permitir la carga, visualización y eliminación de imágenes asociadas a los productos del catálogo, mejorando la identificación visual en el punto de venta y en los reportes.} \\ \hline
		\rowcolor{tableheaderblue}
		\multicolumn{4}{|l|}{\textbf{Funcionalidad}}                                                                                                                                                   \\ \hline
		\multicolumn{4}{|p{0.96\linewidth}|}{%
		El sistema debe permitir cargar imágenes de productos en formatos JPG, PNG o WEBP con un tamaño máximo de 5MB. \newline
		El sistema debe optimizar automáticamente las imágenes cargadas para reducir el tiempo de carga sin pérdida significativa de calidad. \newline
		El sistema debe almacenar las imágenes en un servicio de almacenamiento en la nube (Cloudinary) y guardar únicamente la URL en la base de datos. \newline
		El sistema debe permitir reemplazar la imagen de un producto existente.}                                                                                                                       \\ \hline
		\cellcolor{tableheaderblue}\textbf{Criterios de aceptación} & \multicolumn{3}{p{0.70\linewidth}|}{                                                                                             %
			1. El sistema debe validar el formato y tamaño del archivo antes de iniciar la carga. \newline
			2. El sistema debe mostrar una vista previa de la imagen antes de confirmar su guardado. \newline
		3. El sistema debe mostrar la imagen del producto en el módulo de punto de venta y en el comprobante de venta.}                                                                                \\ \hline
		\rowcolor{tableheaderblue}
		\multicolumn{4}{|l|}{\textbf{Restricciones}}                                                                                                                                                   \\ \hline
		\multicolumn{4}{|p{0.96\linewidth}|}{%
		El sistema debe implementar mecanismos de seguridad para prevenir la carga de archivos maliciosos disfrazados como imágenes.}                                                                  \\ \hline
	
\end{longtable}
\endgroup


\begingroup
\renewcommand{\arraystretch}{1.3}
\setlength{\tabcolsep}{6pt}
\centering
\begin{longtable}{|p{0.25\linewidth}|p{0.15\linewidth}|p{0.15\linewidth}|p{0.35\linewidth}|}
\caption{RF15 - BÚSQUEDA Y FILTRADO AVANZADO DE PRODUCTOS}\label{tab:rf15}\\
\hline
		\rowcolor{tableheaderblue}
		\multicolumn{4}{|l|}{\textbf{Requerimientos funcionales}}                                                                                                                                                  \\ \hline
		\multicolumn{4}{|l|}{\textbf{Prioridad: Alta}}                                                                                                                                                             \\ \hline
		\textbf{CÓDIGO DEL REQUERIMIENTO:}                          & RF15                                                                    & \cellcolor{tableheaderblue}\textbf{Actor} & Administrador / Cajero \\ \hline
		\textbf{NOMBRE DEL REQUERIMIENTO}                           & \multicolumn{3}{l|}{\textbf{Búsqueda y filtrado avanzado de productos}}                                                                      \\ \hline
		\rowcolor{tableheaderblue}
		\multicolumn{4}{|l|}{\textbf{Descripción}}                                                                                                                                                                 \\ \hline
		\multicolumn{4}{|p{0.96\linewidth}|}{%
		El sistema debe proporcionar capacidades de búsqueda y filtrado avanzadas para localizar productos rápidamente en el catálogo mediante múltiples criterios combinables.}                                   \\ \hline
		\rowcolor{tableheaderblue}
		\multicolumn{4}{|l|}{\textbf{Funcionalidad}}                                                                                                                                                               \\ \hline
		\multicolumn{4}{|p{0.96\linewidth}|}{%
		El sistema debe permitir buscar productos por nombre, SKU, código de barras o categoría de forma individual o combinada. \newline
		El sistema debe implementar búsqueda incremental que muestre resultados mientras el usuario escribe. \newline
		El sistema debe permitir filtrar productos por rango de precios, estado (activo/inactivo) y nivel de stock (normal, bajo, agotado). \newline
		El sistema debe ordenar los resultados por nombre, precio, stock o fecha de creación en orden ascendente o descendente.}                                                                                   \\ \hline
		\cellcolor{tableheaderblue}\textbf{Criterios de aceptación} & \multicolumn{3}{p{0.70\linewidth}|}{                                                                                                         %
			1. El sistema debe mostrar resultados de búsqueda en menos de 2 segundos para catálogos de hasta 10,000 productos. \newline
			2. El sistema debe resaltar visualmente los términos de búsqueda coincidentes en los resultados mostrados. \newline
		3. El sistema debe mantener activos los filtros aplicados al navegar entre páginas de resultados.}                                                                                                         \\ \hline
		\rowcolor{tableheaderblue}
		\multicolumn{4}{|l|}{\textbf{Restricciones}}                                                                                                                                                               \\ \hline
		\multicolumn{4}{|p{0.96\linewidth}|}{%
		La búsqueda debe ser insensible a mayúsculas, minúsculas y acentos para mejorar la experiencia del usuario.}                                                                                               \\ \hline
	
\end{longtable}
\endgroup


\begingroup
\renewcommand{\arraystretch}{1.3}
\setlength{\tabcolsep}{6pt}
\centering
\begin{longtable}{|p{0.25\linewidth}|p{0.15\linewidth}|p{0.15\linewidth}|p{0.35\linewidth}|}
\caption{RF16 - APLICACIÓN DE DESCUENTOS EN VENTAS}\label{tab:rf16}\\
\hline
		\rowcolor{tableheaderblue}
		\multicolumn{4}{|l|}{\textbf{Requerimientos funcionales}}                                                                                                                                           \\ \hline
		\multicolumn{4}{|l|}{\textbf{Prioridad: Alta}}                                                                                                                                                      \\ \hline
		\textbf{CÓDIGO DEL REQUERIMIENTO:}                          & RF16                                                             & \cellcolor{tableheaderblue}\textbf{Actor} & Cajero / Administrador \\ \hline
		\textbf{NOMBRE DEL REQUERIMIENTO}                           & \multicolumn{3}{l|}{\textbf{Aplicación de descuentos en ventas}}                                                                      \\ \hline
		\rowcolor{tableheaderblue}
		\multicolumn{4}{|l|}{\textbf{Descripción}}                                                                                                                                                          \\ \hline
		\multicolumn{4}{|p{0.96\linewidth}|}{%
		El sistema debe permitir aplicar descuentos porcentuales o en valor fijo sobre items individuales o sobre el total de la venta, con el objetivo de facilitar promociones y ajustes comerciales.}    \\ \hline
		\rowcolor{tableheaderblue}
		\multicolumn{4}{|l|}{\textbf{Funcionalidad}}                                                                                                                                                        \\ \hline
		\multicolumn{4}{|p{0.96\linewidth}|}{%
		El sistema debe permitir aplicar descuentos de tipo porcentual (\%) o valor fijo (\$) a items individuales en el carrito de venta. \newline
		El sistema debe permitir aplicar un descuento general sobre el subtotal de la venta completa. \newline
		El sistema debe recalcular automáticamente el subtotal, impuestos y total de la venta al aplicar o modificar descuentos. \newline
		El sistema debe registrar en el detalle de la venta los descuentos aplicados con su tipo, monto y justificación.}                                                                                   \\ \hline
		\cellcolor{tableheaderblue}\textbf{Criterios de aceptación} & \multicolumn{3}{p{0.70\linewidth}|}{                                                                                                  %
			1. El sistema debe validar que los descuentos porcentuales se encuentren entre 0\% y 100\%. \newline
			2. El sistema debe impedir aplicar descuentos que resulten en un total de venta negativo. \newline
		3. El sistema debe mostrar claramente en el comprobante de venta los descuentos aplicados y el ahorro obtenido por el cliente.}                                                                     \\ \hline
		\rowcolor{tableheaderblue}
		\multicolumn{4}{|l|}{\textbf{Restricciones}}                                                                                                                                                        \\ \hline
		\multicolumn{4}{|p{0.96\linewidth}|}{%
		El sistema debe requerir autorización mediante contraseña de administrador para descuentos superiores al 30\% del valor del item o de la venta total.}                                              \\ \hline
	
\end{longtable}
\endgroup


\begingroup
\renewcommand{\arraystretch}{1.3}
\setlength{\tabcolsep}{6pt}
\centering
\begin{longtable}{|p{0.25\linewidth}|p{0.15\linewidth}|p{0.15\linewidth}|p{0.35\linewidth}|}
\caption{RF17 - GESTIÓN DEL CARRITO DE COMPRAS TEMPORAL}\label{tab:rf17}\\
\hline
		\rowcolor{tableheaderblue}
		\multicolumn{4}{|l|}{\textbf{Requerimientos funcionales}}                                                                                                                                \\ \hline
		\multicolumn{4}{|l|}{\textbf{Prioridad: Alta}}                                                                                                                                           \\ \hline
		\textbf{CÓDIGO DEL REQUERIMIENTO:}                          & RF17                                                                  & \cellcolor{tableheaderblue}\textbf{Actor} & Cajero \\ \hline
		\textbf{NOMBRE DEL REQUERIMIENTO}                           & \multicolumn{3}{l|}{\textbf{Gestión del carrito de compras temporal}}                                                      \\ \hline
		\rowcolor{tableheaderblue}
		\multicolumn{4}{|l|}{\textbf{Descripción}}                                                                                                                                               \\ \hline
		\multicolumn{4}{|p{0.96\linewidth}|}{%
		El sistema debe permitir gestionar items en un carrito de compras antes de confirmar la venta, facilitando modificaciones de cantidad, eliminación de productos y cálculos dinámicos.}   \\ \hline
		\rowcolor{tableheaderblue}
		\multicolumn{4}{|l|}{\textbf{Funcionalidad}}                                                                                                                                             \\ \hline
		\multicolumn{4}{|p{0.96\linewidth}|}{%
		El sistema debe permitir agregar productos al carrito especificando la cantidad deseada. \newline
		El sistema debe permitir modificar la cantidad de un item ya agregado al carrito mediante incremento, decremento o ingreso directo del valor. \newline
		El sistema debe permitir eliminar items individuales del carrito antes de confirmar la venta. \newline
		El sistema debe permitir vaciar completamente el carrito para iniciar una nueva transacción. \newline
		El sistema debe mostrar en tiempo real el subtotal por item, el subtotal general, impuestos y total de la venta.}                                                                        \\ \hline
		\cellcolor{tableheaderblue}\textbf{Criterios de aceptación} & \multicolumn{3}{p{0.70\linewidth}|}{                                                                                       %
			1. El sistema debe validar disponibilidad de stock al agregar o modificar la cantidad de un producto en el carrito. \newline
			2. El sistema debe solicitar confirmación antes de vaciar completamente el carrito. \newline
		3. El sistema debe mantener el carrito activo durante la sesión del usuario hasta que se confirme la venta o se cierre sesión.}                                                          \\ \hline
		\rowcolor{tableheaderblue}
		\multicolumn{4}{|l|}{\textbf{Restricciones}}                                                                                                                                             \\ \hline
		\multicolumn{4}{|p{0.96\linewidth}|}{%
		El sistema debe limitar la cantidad máxima por item al stock disponible actual del producto.}                                                                                            \\ \hline
	
\end{longtable}
\endgroup


\begingroup
\renewcommand{\arraystretch}{1.3}
\setlength{\tabcolsep}{6pt}
\centering
\begin{longtable}{|p{0.25\linewidth}|p{0.15\linewidth}|p{0.15\linewidth}|p{0.35\linewidth}|}
\caption{RF18 - CONSULTA DE VENTAS POR RANGO DE FECHAS}\label{tab:rf18}\\
\hline
		\rowcolor{tableheaderblue}
		\multicolumn{4}{|l|}{\textbf{Requerimientos funcionales}}                                                                                                                                      \\ \hline
		\multicolumn{4}{|l|}{\textbf{Prioridad: Media}}                                                                                                                                                \\ \hline
		\textbf{CÓDIGO DEL REQUERIMIENTO:}                          & RF18                                                                 & \cellcolor{tableheaderblue}\textbf{Actor} & Administrador \\ \hline
		\textbf{NOMBRE DEL REQUERIMIENTO}                           & \multicolumn{3}{l|}{\textbf{Consulta de ventas por rango de fechas}}                                                             \\ \hline
		\rowcolor{tableheaderblue}
		\multicolumn{4}{|l|}{\textbf{Descripción}}                                                                                                                                                     \\ \hline
		\multicolumn{4}{|p{0.96\linewidth}|}{%
		El sistema debe permitir consultar el historial de ventas realizadas filtrando por un rango de fechas específico, facilitando análisis de períodos comerciales y auditorías.}                  \\ \hline
		\rowcolor{tableheaderblue}
		\multicolumn{4}{|l|}{\textbf{Funcionalidad}}                                                                                                                                                   \\ \hline
		\multicolumn{4}{|p{0.96\linewidth}|}{%
		El sistema debe permitir seleccionar una fecha de inicio y una fecha de fin para delimitar el período de consulta. \newline
		El sistema debe mostrar todas las ventas registradas dentro del rango especificado con su número, fecha, cliente, total y estado. \newline
		El sistema debe calcular y mostrar totales acumulados de ventas, cantidad de transacciones y ticket promedio para el período consultado. \newline
		El sistema debe permitir ordenar los resultados por fecha, número de venta o monto total.}                                                                                                     \\ \hline
		\cellcolor{tableheaderblue}\textbf{Criterios de aceptación} & \multicolumn{3}{p{0.70\linewidth}|}{                                                                                             %
			1. El sistema debe validar que la fecha de inicio no sea posterior a la fecha de fin. \newline
			2. El sistema debe permitir seleccionar rangos predefinidos (hoy, últimos 7 días, último mes, mes actual). \newline
		3. El sistema debe permitir exportar los resultados de la consulta en formato Excel o PDF.}                                                                                                    \\ \hline
		\rowcolor{tableheaderblue}
		\multicolumn{4}{|l|}{\textbf{Restricciones}}                                                                                                                                                   \\ \hline
		\multicolumn{4}{|p{0.96\linewidth}|}{%
		Las consultas de rangos superiores a un año deben implementar paginación para optimizar el rendimiento del sistema.}                                                                           \\ \hline
	
\end{longtable}
\endgroup


\begingroup
\renewcommand{\arraystretch}{1.3}
\setlength{\tabcolsep}{6pt}
\centering
\begin{longtable}{|p{0.25\linewidth}|p{0.15\linewidth}|p{0.15\linewidth}|p{0.35\linewidth}|}
\caption{RF19 - REGISTRO DE MÚLTIPLES CONTACTOS POR CLIENTE}\label{tab:rf19}\\
\hline
		\rowcolor{tableheaderblue}
		\multicolumn{4}{|l|}{\textbf{Requerimientos funcionales}}                                                                                                                                                    \\ \hline
		\multicolumn{4}{|l|}{\textbf{Prioridad: Baja}}                                                                                                                                                               \\ \hline
		\textbf{CÓDIGO DEL REQUERIMIENTO:}                          & RF19                                                                      & \cellcolor{tableheaderblue}\textbf{Actor} & Administrador / Cajero \\ \hline
		\textbf{NOMBRE DEL REQUERIMIENTO}                           & \multicolumn{3}{l|}{\textbf{Registro de múltiples contactos por cliente}}                                                                      \\ \hline
		\rowcolor{tableheaderblue}
		\multicolumn{4}{|l|}{\textbf{Descripción}}                                                                                                                                                                   \\ \hline
		\multicolumn{4}{|p{0.96\linewidth}|}{%
		El sistema debe permitir registrar múltiples números telefónicos y direcciones de correo electrónico para un mismo cliente, facilitando la comunicación por diferentes canales.}                             \\ \hline
		\rowcolor{tableheaderblue}
		\multicolumn{4}{|l|}{\textbf{Funcionalidad}}                                                                                                                                                                 \\ \hline
		\multicolumn{4}{|p{0.96\linewidth}|}{%
		El sistema debe permitir agregar múltiples teléfonos para un cliente, clasificándolos por tipo (móvil, fijo, trabajo). \newline
		El sistema debe permitir agregar múltiples correos electrónicos para un cliente. \newline
		El sistema debe permitir marcar un teléfono y un correo como contacto principal. \newline
		El sistema debe validar el formato de teléfonos y correos electrónicos antes de guardarlos.}                                                                                                                 \\ \hline
		\cellcolor{tableheaderblue}\textbf{Criterios de aceptación} & \multicolumn{3}{p{0.70\linewidth}|}{                                                                                                           %
			1. El sistema debe permitir registrar al menos 3 teléfonos y 3 correos por cliente. \newline
			2. El sistema debe impedir guardar teléfonos o correos duplicados para el mismo cliente. \newline
		3. El sistema debe mostrar los contactos principales destacados visualmente en la ficha del cliente.}                                                                                                        \\ \hline
		\rowcolor{tableheaderblue}
		\multicolumn{4}{|l|}{\textbf{Restricciones}}                                                                                                                                                                 \\ \hline
		\multicolumn{4}{|p{0.96\linewidth}|}{%
		El sistema debe requerir al menos un método de contacto (teléfono o correo) para poder registrar un cliente.}                                                                                                \\ \hline
	
\end{longtable}
\endgroup


\begingroup
\renewcommand{\arraystretch}{1.3}
\setlength{\tabcolsep}{6pt}
\centering
\begin{longtable}{|p{0.25\linewidth}|p{0.15\linewidth}|p{0.15\linewidth}|p{0.35\linewidth}|}
\caption{RF20 - SEGMENTACIÓN DE CLIENTES POR CRITERIOS}\label{tab:rf20}\\
\hline
		\rowcolor{tableheaderblue}
		\multicolumn{4}{|l|}{\textbf{Requerimientos funcionales}}                                                                                                                                           \\ \hline
		\multicolumn{4}{|l|}{\textbf{Prioridad: Baja}}                                                                                                                                                      \\ \hline
		\textbf{CÓDIGO DEL REQUERIMIENTO:}                          & RF20                                                                 & \cellcolor{tableheaderblue}\textbf{Actor} & Administrador      \\ \hline
		\textbf{NOMBRE DEL REQUERIMIENTO}                           & \multicolumn{3}{l|}{\textbf{Segmentación de clientes por criterios}}                                                                  \\ \hline
		\rowcolor{tableheaderblue}
		\multicolumn{4}{|l|}{\textbf{Descripción}}                                                                                                                                                          \\ \hline
		\multicolumn{4}{|p{0.96\linewidth}|}{%
		El sistema debe permitir clasificar y segmentar clientes según criterios comerciales como frecuencia de compra, monto acumulado o fecha de última compra, facilitando estrategias de fidelización.} \\ \hline
		\rowcolor{tableheaderblue}
		\multicolumn{4}{|l|}{\textbf{Funcionalidad}}                                                                                                                                                        \\ \hline
		\multicolumn{4}{|p{0.96\linewidth}|}{%
		El sistema debe calcular automáticamente métricas de comportamiento de compra para cada cliente (total gastado, número de compras, fecha de última compra, ticket promedio). \newline
		El sistema debe permitir crear segmentos personalizados de clientes basados en rangos de las métricas calculadas. \newline
		El sistema debe clasificar automáticamente clientes en categorías predefinidas (VIP, frecuente, ocasional, inactivo) según reglas configurables. \newline
		El sistema debe permitir generar listados de clientes filtrados por segmento o categoría.}                                                                                                          \\ \hline
		\cellcolor{tableheaderblue}\textbf{Criterios de aceptación} & \multicolumn{3}{p{0.70\linewidth}|}{                                                                                                  %
			1. El sistema debe recalcular las métricas de comportamiento cada vez que se registre una nueva venta asociada al cliente. \newline
			2. El sistema debe permitir al administrador configurar los umbrales que definen cada categoría de cliente. \newline
		3. El sistema debe mostrar visualmente la categoría del cliente en su ficha y en el punto de venta.}                                                                                                \\ \hline
		\rowcolor{tableheaderblue}
		\multicolumn{4}{|l|}{\textbf{Restricciones}}                                                                                                                                                        \\ \hline
		\multicolumn{4}{|p{0.96\linewidth}|}{%
		Los segmentos y categorías deben recalcularse mediante procesos programados para no afectar el rendimiento del sistema durante horarios de operación.}                                              \\ \hline
	
\end{longtable}
\endgroup


\begingroup
\renewcommand{\arraystretch}{1.3}
\setlength{\tabcolsep}{6pt}
\centering
\begin{longtable}{|p{0.25\linewidth}|p{0.15\linewidth}|p{0.15\linewidth}|p{0.35\linewidth}|}
\caption{RF21 - REPORTE DE PRODUCTOS MÁS VENDIDOS}\label{tab:rf21}\\
\hline
		\rowcolor{tableheaderblue}
		\multicolumn{4}{|l|}{\textbf{Requerimientos funcionales}}                                                                                                                                 \\ \hline
		\multicolumn{4}{|l|}{\textbf{Prioridad: Media}}                                                                                                                                           \\ \hline
		\textbf{CÓDIGO DEL REQUERIMIENTO:}                          & RF21                                                            & \cellcolor{tableheaderblue}\textbf{Actor} & Administrador \\ \hline
		\textbf{NOMBRE DEL REQUERIMIENTO}                           & \multicolumn{3}{l|}{\textbf{Reporte de productos más vendidos}}                                                             \\ \hline
		\rowcolor{tableheaderblue}
		\multicolumn{4}{|l|}{\textbf{Descripción}}                                                                                                                                                \\ \hline
		\multicolumn{4}{|p{0.96\linewidth}|}{%
		El sistema debe generar un reporte que identifique los productos con mayor volumen de ventas en un período determinado, facilitando decisiones de inventario y estrategias comerciales.}  \\ \hline
		\rowcolor{tableheaderblue}
		\multicolumn{4}{|l|}{\textbf{Funcionalidad}}                                                                                                                                              \\ \hline
		\multicolumn{4}{|p{0.96\linewidth}|}{%
		El sistema debe permitir seleccionar un rango de fechas para generar el reporte de productos más vendidos. \newline
		El sistema debe calcular para cada producto la cantidad total vendida y el valor total generado en el período. \newline
		El sistema debe ordenar los resultados por cantidad vendida o por valor total de forma descendente. \newline
		El sistema debe permitir limitar el reporte a un número específico de productos (top 10, top 20, top 50). \newline
		El sistema debe presentar los resultados en formato de tabla y gráfica de barras.}                                                                                                        \\ \hline
		\cellcolor{tableheaderblue}\textbf{Criterios de aceptación} & \multicolumn{3}{p{0.70\linewidth}|}{                                                                                        %
			1. El sistema debe excluir del reporte las ventas anuladas o devueltas. \newline
			2. El sistema debe permitir filtrar el reporte por categoría de productos. \newline
		3. El sistema debe permitir exportar el reporte en formato Excel o PDF.}                                                                                                                  \\ \hline
		\rowcolor{tableheaderblue}
		\multicolumn{4}{|l|}{\textbf{Restricciones}}                                                                                                                                              \\ \hline
		\multicolumn{4}{|p{0.96\linewidth}|}{%
		El reporte debe considerar únicamente ventas con estado "completada" para garantizar precisión de la información.}                                                                        \\ \hline
	
\end{longtable}
\endgroup


\begingroup
\renewcommand{\arraystretch}{1.3}
\setlength{\tabcolsep}{6pt}
\centering
\begin{longtable}{|p{0.25\linewidth}|p{0.15\linewidth}|p{0.15\linewidth}|p{0.35\linewidth}|}
\caption{RF22 - ANÁLISIS DE RENTABILIDAD POR PRODUCTO}\label{tab:rf22}\\
\hline
		\rowcolor{tableheaderblue}
		\multicolumn{4}{|l|}{\textbf{Requerimientos funcionales}}                                                                                                                                     \\ \hline
		\multicolumn{4}{|l|}{\textbf{Prioridad: Media}}                                                                                                                                               \\ \hline
		\textbf{CÓDIGO DEL REQUERIMIENTO:}                          & RF22                                                                & \cellcolor{tableheaderblue}\textbf{Actor} & Administrador \\ \hline
		\textbf{NOMBRE DEL REQUERIMIENTO}                           & \multicolumn{3}{l|}{\textbf{Análisis de rentabilidad por producto}}                                                             \\ \hline
		\rowcolor{tableheaderblue}
		\multicolumn{4}{|l|}{\textbf{Descripción}}                                                                                                                                                    \\ \hline
		\multicolumn{4}{|p{0.96\linewidth}|}{%
		El sistema debe calcular y presentar métricas de rentabilidad por producto comparando precios de costo y venta, facilitando la identificación de productos con mejores márgenes comerciales.} \\ \hline
		\rowcolor{tableheaderblue}
		\multicolumn{4}{|l|}{\textbf{Funcionalidad}}                                                                                                                                                  \\ \hline
		\multicolumn{4}{|p{0.96\linewidth}|}{%
		El sistema debe calcular el margen bruto unitario por producto (precio de venta - precio de costo). \newline
		El sistema debe calcular el margen bruto porcentual por producto ((precio de venta - precio de costo) / precio de venta * 100). \newline
		El sistema debe calcular la rentabilidad total acumulada por producto considerando las cantidades vendidas en un período. \newline
		El sistema debe permitir ordenar productos por margen porcentual o rentabilidad total. \newline
		El sistema debe identificar visualmente productos con márgenes negativos o muy bajos (alertas).}                                                                                              \\ \hline
		\cellcolor{tableheaderblue}\textbf{Criterios de aceptación} & \multicolumn{3}{p{0.70\linewidth}|}{                                                                                            %
			1. El sistema debe actualizar los cálculos de rentabilidad cuando se modifiquen los precios de costo o venta de un producto. \newline
			2. El sistema debe permitir filtrar el análisis por categoría y rango de fechas de ventas. \newline
		3. El sistema debe presentar un resumen con el margen promedio general del negocio.}                                                                                                          \\ \hline
		\rowcolor{tableheaderblue}
		\multicolumn{4}{|l|}{\textbf{Restricciones}}                                                                                                                                                  \\ \hline
		\multicolumn{4}{|p{0.96\linewidth}|}{%
		Los cálculos de rentabilidad deben considerar el precio de costo y venta vigentes al momento de cada transacción registrada.}                                                                 \\ \hline
	
\end{longtable}
\endgroup


\begingroup
\renewcommand{\arraystretch}{1.3}
\setlength{\tabcolsep}{6pt}
\centering
\begin{longtable}{|p{0.25\linewidth}|p{0.15\linewidth}|p{0.15\linewidth}|p{0.35\linewidth}|}
\caption{RF23 - REPORTE DE VENTAS POR MÉTODO DE PAGO}\label{tab:rf23}\\
\hline
		\rowcolor{tableheaderblue}
		\multicolumn{4}{|l|}{\textbf{Requerimientos funcionales}}                                                                                                                                    \\ \hline
		\multicolumn{4}{|l|}{\textbf{Prioridad: Baja}}                                                                                                                                               \\ \hline
		\textbf{CÓDIGO DEL REQUERIMIENTO:}                          & RF23                                                               & \cellcolor{tableheaderblue}\textbf{Actor} & Administrador \\ \hline
		\textbf{NOMBRE DEL REQUERIMIENTO}                           & \multicolumn{3}{l|}{\textbf{Reporte de ventas por método de pago}}                                                             \\ \hline
		\rowcolor{tableheaderblue}
		\multicolumn{4}{|l|}{\textbf{Descripción}}                                                                                                                                                   \\ \hline
		\multicolumn{4}{|p{0.96\linewidth}|}{%
		El sistema debe generar un reporte que agrupe las ventas según el método de pago utilizado, facilitando la conciliación financiera y el análisis de preferencias de pago de los clientes.}   \\ \hline
		\rowcolor{tableheaderblue}
		\multicolumn{4}{|l|}{\textbf{Funcionalidad}}                                                                                                                                                 \\ \hline
		\multicolumn{4}{|p{0.96\linewidth}|}{%
		El sistema debe permitir seleccionar un rango de fechas para generar el reporte de ventas por método de pago. \newline
		El sistema debe calcular el total de ventas y cantidad de transacciones para cada método de pago (efectivo, tarjeta, transferencia). \newline
		El sistema debe calcular el porcentaje que representa cada método de pago sobre el total de ventas del período. \newline
		El sistema debe presentar los resultados en formato de tabla y gráfica circular (pie chart).}                                                                                                \\ \hline
		\cellcolor{tableheaderblue}\textbf{Criterios de aceptación} & \multicolumn{3}{p{0.70\linewidth}|}{                                                                                           %
			1. El sistema debe excluir del reporte las ventas anuladas. \newline
			2. El sistema debe mostrar el método de pago más utilizado destacado visualmente. \newline
		3. El sistema debe permitir exportar el reporte en formato PDF o Excel.}                                                                                                                     \\ \hline
		\rowcolor{tableheaderblue}
		\multicolumn{4}{|l|}{\textbf{Restricciones}}                                                                                                                                                 \\ \hline
		\multicolumn{4}{|p{0.96\linewidth}|}{%
		El reporte debe considerar únicamente ventas confirmadas y no debe incluir transacciones pendientes o en proceso.}                                                                           \\ \hline
	
\end{longtable}
\endgroup


\begingroup
\renewcommand{\arraystretch}{1.3}
\setlength{\tabcolsep}{6pt}
\centering
\begin{longtable}{|p{0.25\linewidth}|p{0.15\linewidth}|p{0.15\linewidth}|p{0.35\linewidth}|}
\caption{RF24 - REPORTE DE MOVIMIENTOS DE INVENTARIO}\label{tab:rf24}\\
\hline
		\rowcolor{tableheaderblue}
		\multicolumn{4}{|l|}{\textbf{Requerimientos funcionales}}                                                                                                                                               \\ \hline
		\multicolumn{4}{|l|}{\textbf{Prioridad: Media}}                                                                                                                                                         \\ \hline
		\textbf{CÓDIGO DEL REQUERIMIENTO:}                          & RF24                                                               & \cellcolor{tableheaderblue}\textbf{Actor} & Administrador            \\ \hline
		\textbf{NOMBRE DEL REQUERIMIENTO}                           & \multicolumn{3}{l|}{\textbf{Reporte de movimientos de inventario}}                                                                        \\ \hline
		\rowcolor{tableheaderblue}
		\multicolumn{4}{|l|}{\textbf{Descripción}}                                                                                                                                                              \\ \hline
		\multicolumn{4}{|p{0.96\linewidth}|}{%
		El sistema debe generar un reporte detallado de todos los movimientos de inventario registrados en un período, clasificados por tipo de movimiento, facilitando la auditoría y trazabilidad del stock.} \\ \hline
		\rowcolor{tableheaderblue}
		\multicolumn{4}{|l|}{\textbf{Funcionalidad}}                                                                                                                                                            \\ \hline
		\multicolumn{4}{|p{0.96\linewidth}|}{%
		El sistema debe permitir seleccionar un rango de fechas para generar el reporte de movimientos de inventario. \newline
		El sistema debe permitir filtrar movimientos por tipo (entrada, salida, ajuste, venta) y por producto específico. \newline
		El sistema debe mostrar para cada movimiento la fecha, tipo, producto, cantidad, stock anterior, stock resultante y usuario responsable. \newline
		El sistema debe calcular totales de entradas y salidas por tipo de movimiento para el período consultado.}                                                                                              \\ \hline
		\cellcolor{tableheaderblue}\textbf{Criterios de aceptación} & \multicolumn{3}{p{0.70\linewidth}|}{                                                                                                      %
			1. El sistema debe ordenar los movimientos cronológicamente del más reciente al más antiguo por defecto. \newline
			2. El sistema debe permitir filtrar movimientos por categoría de producto. \newline
		3. El sistema debe permitir exportar el reporte completo en formato Excel incluyendo todos los detalles.}                                                                                               \\ \hline
		\rowcolor{tableheaderblue}
		\multicolumn{4}{|l|}{\textbf{Restricciones}}                                                                                                                                                            \\ \hline
		\multicolumn{4}{|p{0.96\linewidth}|}{%
		El reporte debe incluir únicamente movimientos confirmados y registrados formalmente en el sistema, excluyendo movimientos en borrador o pendientes de aprobación.}                                     \\ \hline
	
\end{longtable}
\endgroup


